\documentclass[bachelor, och, pract, times]{SCWorks}
% Тип обучения (одно из значений):
%   bachelor   - бакалавриат (по умолчанию)
%   spec       - специальность
%   master     - магистратура
% Форма обучения (одно из значений):
%   och        - очное (по умолчанию)
%   zaoch      - заочное
% Тип работы (одно из значений):
%   coursework - курсовая работа (по умолчанию)
%   referat    - реферат
% * otchet     - универсальный отчет
% * nirjournal - журнал НИР
% * digital    - итоговая работа для цифровой кафедры
%   diploma    - дипломная работа
%   pract      - отчет о научно-исследовательской работе
%   autoref    - автореферат выпускной работы
%   assignment - задание на выпускную квалификационную работу
%   review     - отзыв руководителя
%   critique   - рецензия на выпускную работу

% * Добавлены вручную. За вопросами к @mchernigin
\usepackage{preamble}

\begin{document}

% Кафедра (в родительном падеже)
\chair{математической кибернетики и компьютерных наук}

% Тема работы
\title{Разработка интерактивного веб"=приложения Mesto}

% Курс
\course{3}

% Группа
\group{351}

% Семестр (только для практики, для остальных типов работ не используется)
\term{3}

% Наименование практики (только для практики, для остальных типов работ не
% используется)
\practtype{научно-исследовательская работа (учебная практика, рассредоточенная)}
% Факультет (в родительном падеже) (по умолчанию "факультета КНиИТ")
% \department{факультета КНиИТ}

% Специальность/направление код - наименование
% \napravlenie{02.03.02 "--- Фундаментальная информатика и информационные технологии}
% \napravlenie{02.03.01 "--- Математическое обеспечение и администрирование информационных систем}
% \napravlenie{09.03.01 "--- Информатика и вычислительная техника}
\napravlenie{09.03.04 "--- Программная инженерия}
% \napravlenie{10.05.01 "--- Компьютерная безопасность}

% Для студентки. Для работы студента следующая команда не нужна.
% \studenttitle{студентки}

% Фамилия, имя, отчество в родительном падеже
\author{Смирнова Егора Ильича}

% Заведующий кафедрой
\chtitle{доцент, к.\,ф.-м.\,н.}
\chname{С.\,В.\,Миронов}

% Руководитель ДПП ПП для цифровой кафедры (перекрывает заведующего кафедры)
% \chpretitle{
%     заведующий кафедрой математических основ информатики и олимпиадного\\
%     программирования на базе МАОУ <<Ф"=Т лицей №1>>
% }
% \chtitle{г. Саратов, к.\,ф.-м.\,н., доцент}
% \chname{Кондратова\, Ю.\,Н.}

% Научный руководитель (для реферата преподаватель проверяющий работу)
\satitle{доцент, к.\,ф.-м.\,н.} %должность, степень, звание
\saname{А.\,С.\,Иванова}

% Руководитель практики от организации (руководитель для цифровой кафедры)
\patitle{доцент, к.\,ф.-м.\,н.}
\paname{А.\,С.\,Иванова}

% Семестр (только для практики, для остальных типов работ не используется)
\term{6}

% Наименование практики (только для практики, для остальных типов работ не
% используется)
\practtype{учебная}

% Продолжительность практики (количество недель) (только для практики, для
% остальных типов работ не используется)
\duration{15}

% Даты начала и окончания практики (только для практики, для остальных типов
% работ не используется)
\practStart{08.02.2025}
\practFinish{22.5.2025}

% Год выполнения отчета
\date{2026}

\maketitle

% Включение нумерации рисунков, формул и таблиц по разделам (по умолчанию -
% нумерация сквозная) (допускается оба вида нумерации)
\secNumbering

\tableofcontents

\intro
В условиях современной цифровой среды к клиентским веб-приложениям
предъявляются высочайшие требования с точки зрения интерактивности и надежности
интерфейса. Статические страницы уступили место динамическим сервисам, в
которых взаимодействие пользователя с данными происходит без перезагрузки
системы. Центральную роль в обеспечении такой динамики играет язык
программирования JavaScript, позволяющий модифицировать структуру документа
непосредственно в процессе работы пользователя.\cite{mdn_javascript_overview}

В рамках учебной практики ставилась задача проектирования и реализации
клиентской части интерактивного веб-сервиса «Mesto». Данный проект представляет
собой визуальную платформу для публикации фотографий мест, управления профилем
и взаимодействия с контентом посредством системы веб-отметок. Важной
особенностью разработки стала организация полного цикла, включающего написание
кода, автоматизированную сборку и развертывание в сетевой среде. Интеграция с
сервером реализована в стиле REST API, где действия пользователя сопровождаются
отправкой запросов и синхронизацией интерфейса.

Для сборки продукта использован программный комплекс Vite, обеспечивающий
оптимизацию пакетов файлов и минификацию скриптов. Автоматизация процессов
интеграции и доставки кодовой базы реализована посредством инструментов GitHub
Actions, что позволяет автоматически компилировать проект и осуществлять
публикацию бандла на платформе GitHub Pages. В приложении также функционирует
система программной валидации форм на основе интерфейса контроля ограничений,
обеспечивающая мониторинг ввода в реальном времени. В рамках индивидуального
задания по варианту №2 внедрен аналитический модуль для сбора и вывода
агрегированной статистики пользовательской активности по клику на логотип
сервиса.

Основной целью учебной практики является получение опыта проектирования,
структурирования и автоматизированного развертывания модульных приложений на
языке JavaScript.

Для достижения цели решались следующие практические задачи:
\begin{itemize}
    \item спроектировать модульную архитектуру программного комплекса с
      разделением слоев сетевого взаимодействия, валидации и управления окнами;
    \item организовать взаимодействие с внешним сервером через интерфейс Fetch
      и настроить централизованную обработку ошибок;
    \item разработать систему программной валидации форм с выводом кастомных
      сообщений и контролем состояния кнопок подтверждения;
    \item реализовать безопасный динамический рендеринг элементов на основе
      инертных HTML-шаблонов;
    \item выполнить индивидуальное задание по проектированию алгоритма
      агрегации данных для формирования статистической сводки об активности
      авторов;
    \item настроить конфигурацию сборщика Vite и автоматизировать процесс
      непрерывного деплоя на хостинг GitHub Pages через сценарии GitHub
      Actions.
\end{itemize}

\section{Теоретическая часть}
\subsection{Управление состояниями отложенных операций и объект Promise}
Выполнение сложных вычислительных задач или процедур взаимодействия с внешними
источниками данных требует использования специальных инструментов контроля.
Когда результат операции не может быть получен мгновенно в текущей строке кода,
возникает необходимость в создании управляемой структуры для отслеживания
жизненного цикла данной операции. Современным стандартом управления такими
потоками информации выступает объект Promise.

Объект Promise представляет собой структурированный контейнер, предназначенный
для хранения итогового результата вычислений или полученных данных. Главное
свойство данного объекта заключается в черком разграничении этапов выполнения
операции через фиксированные статусы.

Жизненный цикл объекта Promise включает три строгих состояния:
\begin{itemize}
    \item Состояние ожидания или статус pending. Данный статус присваивается
      объекту немедленно в момент его инициализации. Операция обработки данных
      уже запущена, однако итоговый результат вычислений на текущем этапе
      остается неизвестным.
    \item Разрешенное состояние или статус fulfilled. Программная задача
      успешно завершена, контейнер наполняется искомыми верифицированными
      данными и фиксирует их для дальнейшего использования.
    \item Отклоненное состояние или статус rejected. В процессе обработки
      данных возник критический сбой, объект прекращает выполнение операции и
      сохраняет подробную информацию о характере возникшей ошибки.
\end{itemize}

После завершения целевой задачи состояние контейнера изменяется ровно один раз.
Переход из статуса pending в статус fulfilled или статус rejected является
необратимым. Дальнейшая модификация внутреннего состояния объекта полностью
блокируется средой выполнения.

Для извлечения сохраненных данных разработчик выстраивает последовательные
цепочки вызовов. Специализированный метод then принимает функцию обработки
успешного результата. Данная функция активируется исключительно после перехода
объекта Promise в разрешенное состояние. Метод catch принимает функцию
перехвата исключений, которая вызывается при переходе объекта в статус
rejected. Построение таких цепочек позволяет структурировать логику программы,
избегать глубокой вложенности функций обратного вызова и сохранять линейный вид
кодовой базы \cite{learn_js_promise, mdn_promise}.

Отдельную роль в архитектуре приложения играет механизм параллельной
координации группы независимых операций. Специализированный метод Promise.all
принимает массив объектов Promise и возвращает один объединенный результат.
Данный результат переходит в состояние fulfilled только после успешного
завершения абсолютно всех задач из исходного массива. Если хотя бы один элемент
массива переходит в статус rejected, весь объединенный результат немедленно
завершается ошибкой. Такой подход необходим при первоначальной инициализации
интерфейса, когда для корректного отображения страницы требуется одновременное
наличие нескольких независимых массивов данных.

\subsection{Сетевое взаимодействие и протоколы передачи данных}
Клиентские веб-приложения функционируют в распределенной информационной среде.
Обмен данными с удаленными серверами происходит через протокол передачи
гипертекста HTTP. Современные серверные интерфейсы проектируются на базе
архитектурного стиля REST. В основе данного стиля лежит абстракция ресурса.
Каждая сущность информационной системы имеет уникальный идентификатор в виде
сетевого адреса.

Взаимодействие клиента и сервера не сохраняет состояние между отдельными
запросами. Отправляемый пакет содержит абсолютно всю информацию для корректной
обработки сервером. Клиент манипулирует ресурсами посредством строгих методов:
\begin{itemize}
    \item Метод GET предназначен исключительно для извлечения информации.
      Данный метод является безопасным и идемпотентным, многократное выполнение
      запроса не приводит к изменению состояния базы данных.
    \item Метод POST используется для добавления новых сущностей. Запрос
      содержит сериализованные данные создаваемого объекта.
    \item Метод PATCH применяется для частичной модификации существующего
      ресурса. Сервер обновляет исключительно переданные поля, оставляя
      остальные данные без изменений.
    \item Метод PUT осуществляет полную замену целевого ресурса предоставленным
      набором данных.
    \item Метод DELETE инициирует безвозвратное уничтожение указанных данных на
      стороне сервера \cite{mdn_http_methods}.
\end{itemize}

Передача сложной информации между клиентом и сервером осуществляется в
универсальном текстовом формате JSON. Перед отправкой запроса клиентские
скрипты сериализуют объекты в специальную строку. При получении ответа от
сервера происходит обратный процесс десериализации строки в полноценные
структуры данных.

Каждый запрос сопровождается набором служебных заголовков. Заголовки передают
важную метаинформацию. Заголовок авторизации содержит уникальный секретный
ключ. Ключ подтверждает личность пользователя и дает право на модификацию
защищенных ресурсов. Заголовок типа контента сообщает серверу формат
передаваемого тела запроса, гарантируя корректную расшифровку отправленных
данных.

\subsection{Инкапсуляция операций обработки данных и интерфейсные контракты}
Взаимодействие между главными скриптами и источником данных требует создания
промежуточного абстрактного слоя. Данный слой полностью изолирует механизмы
получения, валидации и преобразования информации от логики рендеринга
графического интерфейса. Создание специализированных функций для выполнения
однотипных операций позволяет унифицировать программный интерфейс и снизить
дублирование кодовой базы.

Каждая функция абстрактного слоя выполняет роль строгого интерфейсного
контракта. Функция принимает на вход определенные параметры, осуществляет их
структурирование, преобразует данные в требуемый формат и возвращает результат
вызывающему компоненту. Вызывающий компонент не обладает информацией о
внутреннем устройстве абстрактного слоя, что позволяет изменять внутреннюю
логику обработки данных без необходимости модификации кода визуального
представления.

Важной частью интерфейсного контракта является унифицированная обработка
исключительных ситуаций. При обнаружении некорректных входных данных или
невозможности выполнить операцию абстрактный слой генерирует ошибку. Ошибка
передается по цепочке вызовов в главный координирующий модуль, где запускаются
заложенные сценарии восстановления работоспособности системы или оповещения
пользователя через графические элементы \cite{mdn_fetch}.

\subsection{Синхронизация состояний графического интерфейса}
Пользовательский интерфейс должен мгновенно отражать изменения, происходящие в
структурах данных приложения. Любое действие человека, приводящее к модификации
информации, инициирует циклическое обновление визуальных элементов страницы.
Процесс синхронизации включает перевод интерактивных компонентов в
соответствующие режимы работы.

Перед началом выполнения транзакции по изменению данных скрипт блокирует
элементы управления для предотвращения повторных кликов. Повторные действия
пользователя могут нарушить целостность структуры данных и вызвать критические
ошибки выполнения кода. Текстовое содержимое управляющего элемента временно
заменяется на служебное сообщение о текущем статусе операции. После успешного
завершения обработки данных и обновления структуры скрипт возвращает
интерфейсные элементы в исходное состояние, обеспечивая согласованность
визуального представления и программного состояния.

\subsection{Изоляция программной логики и модульная архитектура}
Исторически клиентские скрипты выполнялись в едином глобальном пространстве
имен. Объявление переменных в одном файле делало их доступными во всех
остальных файлах проекта. Подобный подход приводил к конфликтам имен и
затруднял масштабирование кодовой базы. Современная парадигма разработки
требует жесткой изоляции программного кода. Изоляция достигается путем
разделения монолитной логики на независимые функциональные блоки.

Каждый блок представляет собой отдельный файл с собственным лексическим
окружением. Внутренние переменные и функции остаются полностью скрытыми от
внешнего кода по умолчанию. Разработчик обязан явно декларировать интерфейс
взаимодействия. Специальные ключевые слова позволяют экспортировать строго
определенные сущности. Другие файлы импортируют необходимые функции для
построения программных связей \cite{mdn_modules}.

Модульная система гарантирует однократное выполнение кода при первом импорте.
Последующие импорты обращаются к кэшированному результату. Данный механизм
экономит вычислительные ресурсы и обеспечивает консистентность данных между
различными частями приложения. Модули автоматически работают в строгом режиме,
блокируя использование небезопасных конструкций языка \cite{morison}.

\subsection{Интерфейс программного контроля данных}
Клиентская валидация данных является обязательным этапом проектирования
надежных интерфейсов. Предварительная фильтрация ввода защищает серверную
инфраструктуру от избыточной нагрузки и некорректных запросов. Одновременно
обеспечивается мгновенная визуальная обратная связь, повышающая качество
взаимодействия пользователя с системой.

Декларативная валидация работает на уровне разметки. Разработчик устанавливает
атрибуты обязательности заполнения, минимальной длины строки и регулярные
выражения. Браузер автоматически блокирует отправку некорректных форм и выводит
системные предупреждения. Визуальное представление системных предупреждений
невозможно изменить средствами каскадных таблиц стилей. Системный подход
серьезно ограничивает гибкость интерфейса.

Профессиональная разработка требует применения программного интерфейса контроля
ограничений. Данный интерфейс предоставляет полный доступ к состоянию каждого
поля ввода. Состояние описывается специальным объектом, содержащим набор
булевых флагов.

Флаг пропуска значения принимает истинное значение при оставлении обязательного
поля пустым. Флаг несоответствия шаблону активируется при нарушении правил
регулярного выражения. Флаг несоответствия типа указывает на ввод недопустимых
символов в специализированные поля ввода. Глобальный флаг валидности
подтверждает абсолютное соответствие введенного текста всем установленным
правилам \cite{mdn_validity_state}.

Скрипт перехватывает событие отправки формы и отменяет стандартное поведение
браузера. При обнаружении ошибок ввода вызывается метод установки
пользовательского предупреждения. Этот метод заменяет системный текст на
сообщение, заранее подготовленное разработчиком в атрибутах разметки.

Архитектура валидации включает непрерывный мониторинг состояния всей формы.
Проверка запускается при каждом изменении содержимого полей ввода. Скрипт
анализирует флаги валидности каждого элемента. Кнопка подтверждения действия
получает атрибут блокировки до полного исправления всех допущенных ошибок.
Блокировка предотвращает отправку заведомо неверных данных и делает интерфейс
предсказуемым.

\subsection{Объектная модель документа и изолированная шаблонизация}
Объектная модель документа представляет структуру веб-страницы в виде
иерархического дерева взаимодействующих узлов. Скрипты используют этот
программный интерфейс для поиска элементов, модификации их содержимого и
реакции на действия пользователя.

Изменение активного дерева узлов является наиболее ресурсоемкой операцией в
браузере. Добавление новых визуальных элементов инициирует сложный процесс
перерасчета размеров всей страницы. За перерасчетом следует процесс
попиксельной перерисовки обновленных областей экрана. Неоптимизированные
модификации приводят к задержкам вывода графики.

Генерация новой разметки путем объединения текстовых строк представляет
критическую угрозу безопасности. Внедрение строк заставляет браузер заново
запускать анализатор разметки. Передача пользовательского текста напрямую в
структуру страницы открывает уязвимость межсайтового скриптинга. Вредоносный
скрипт получает возможность исполниться под видом обычного текста.

Современный стандарт безопасного рендеринга базируется на использовании
специализированного тега шаблона. Данный тег обладает свойством инертности.
Браузер парсит содержимое шаблона при первоначальной загрузке документа, но
полностью исключает его из процесса отрисовки. Вложенные медиафайлы не
скачиваются, а прописанные стили не применяются \cite{mdn_template}.

\subsection{Инструменты автоматизации сборки и оптимизации распределения ресурсов}
Разработка современных клиентских приложений требует применения
специализированных программных комплексов для подготовки исходного кода к
публикации в глобальной сети. Инструменты автоматизации сборки осуществляют
глубокий статический анализ кодовой базы и формируют оптимизированный пакет
файлов. Процесс подготовки включает компиляцию разрозненных независимых модулей
в единую связную структуру. На данном этапе применяются алгоритмы минификации и
обфускации. Указанные алгоритмы удаляют служебные пробелы, вырезают комментарии
разработчиков и преобразуют имена переменных в короткие нечитаемые
последовательности, что радикально уменьшает итоговый размер скриптов и
ускоряет их скачивание клиентом.

Отдельным этапом сборки выступает настройка путей распределения статических
ресурсов. Для корректного функционирования графического интерфейса на удаленном
сервере инструмент сборки автоматически трансформирует абсолютные локальные
пути файловой системы в безопасные относительные ссылки. Подобная трансформация
гарантирует правильную загрузку изображений, шрифтов и каскадных таблиц стилей
независимо от внутренней иерархии каталогов на конечном веб-сервере.\cite{vite_docs}.

\subsection{Непрерывная интеграция и автоматизированное развертывание в сетевой среде}
Публикация готового программного продукта базируется на методологии непрерывной
интеграции и автоматизированной доставки кода. Данная методология полностью
исключает необходимость ручного переноса файлов на сервер и сводит к нулю
вероятность возникновения человеческой ошибки при обновлении системы.
Управление жизненным циклом продукта осуществляется через платформы
автоматизации рабочих процессов, которые используют декларативные
конфигурационные сценарии. Сценарий описывает строгую последовательность
команд, выполняемых изолированным виртуальным окружением.

Триггером для запуска автоматизированного сценария выступает любое обновление
исходного кода в главной ветке удаленного репозитория. Обнаружив изменения,
платформа инициализирует среду выполнения, скачивает необходимые зависимости и
запускает процесс компиляции проекта. После успешного завершения сборки скрипт
изолирует директорию с оптимизированными файлами и инициирует процесс
публикации на платформе статического хостинга. Синхронизация файлов происходит
в фоновом режиме, мгновенно делая самую актуальную версию интерфейса доступной
для конечных пользователей.

\section{Практическая часть}
\subsection{Файловая структура и архитектура приложения}
Проект собирается инструментом Vite и организован по принципу строгой модульной
архитектуры. Стили оформлены по методологии БЭМ и хранятся в директории
исходных кодов в виде отдельных файлов для каждого визуального блока. Скрипты
разделены на компонентные модули.

Разработанная файловая структура включает следующие элементы:
\begin{itemize}
    \item файл \texttt{index.html} — корневой документ с семантической
      разметкой, всплывающими окнами и скрытыми шаблонами генерации карточек и
      элементов статистики;
    \item файл \texttt{src/scripts/index.js} — главная точка входа:
      инициализация ссылок на узлы документа, регистрация слушателей событий,
      запуск контроля ввода и первичная загрузка данных;
    \item файл \texttt{src/scripts/components/api.js} — сетевой слой:
      конфигурация подключения и все запросы протокола передачи гипертекста;
    \item файл \texttt{src/scripts/components/card.js} — компонент генерации
      карточки: клонирование узла из шаблона, управление отметками лайков и
      удалением элементов;
    \item файл \texttt{src/scripts/components/modal.js} — модуль управления
      состояниями видимости модальных окон;
    \item файл \texttt{src/scripts/components/validation.js} — компонент
      программного контроля пользовательского ввода;
    \item файл \texttt{vite.config.js} — конфигурация сборщика для оптимизации
      путей распределения ресурсов;
    \item файл \texttt{.github/workflows/deploy.yml} — конфигурационный
      сценарий автоматического развертывания проекта.
\end{itemize}

Директория скомпилированных файлов \texttt{dist} формируется командой сборки и публикуется во внешний публичный репозиторий платформы GitHub Pages посредством автоматизированных сценариев GitHub Actions.

\subsection{Модуль сетевых запросов}
Весь сетевой слой сосредоточен в изолированном файле. Конфигурация запросов
вынесена в отдельный объект, содержащий базовый сетевой адрес и служебные
заголовки.

\begin{minted}{javascript}
const config = {
  baseUrl: "https://mesto.nomoreparties.co/v1/{id}",
  headers: {
    authorization: "{token}",
    "Content-Type": "application/json",
  },
};
\end{minted}

Функция обработки ответа служит единой точкой проверки статуса выполнения
операции. Данная функция проверяет свойство успешности и при ошибочном статусе
принудительно отклоняет объект Promise.

\begin{minted}{javascript}
const getResponseData = (res) => {
  return res.ok ? res.json() : Promise.reject(`Ошибка: ${res.status}`);
};
\end{minted}

Каждая экспортируемая функция инкапсулирует один запрос. Функции обновления
профиля и ссылки на изображение аватара используют метод PATCH и передают тело
запроса в сериализованном формате JSON. Для переключения лайка применяется
единая функция, динамически выбирающая метод взаимодействия. Функция отправляет
метод DELETE при снятии отметки и метод PUT при ее установке.

\begin{minted}{javascript}
export const changeLikeCardStatus = (cardId, isLiked) => {
  return fetch(`${config.baseUrl}/cards/likes/${cardId}`, {
    method: isLiked ? "DELETE" : "PUT",
    headers: config.headers,
  }).then(getResponseData);
};
\end{minted}

\subsection{Компонент управления карточками}
Модуль управления карточками реализует создание визуального элемента и контроль
интерактивных действий. Извлечение шаблона вынесено в отдельную локальную
функцию, что упрощает чтение кода и обеспечивает изоляцию манипуляций с деревом
документа.

Главная функция генерации принимает данные карточки, объект функций обратного
вызова и идентификатор текущего пользователя. Кнопка удаления извлекается из
разметки и физически удаляется из дерева узлов для карточек, созданных другими
участниками системы.

\begin{minted}{javascript}
export const createCardElement = (
  data,
  { onPreviewPicture, onLikeIcon, onDeleteCard },
  currentUserId
) => {
  const cardElement = getTemplate();
  const likeButton = cardElement.querySelector(".card__like-button");
  const deleteButton = cardElement.querySelector(".card__control-button_type_delete");
  const cardImage = cardElement.querySelector(".card__image");
  const likeCountElement = cardElement.querySelector(".card__like-count");
  cardImage.src = data.link;
  cardImage.alt = data.name;
  cardElement.querySelector(".card__title").textContent = data.name;
  likeCountElement.textContent = data.likes.length;
  const isLikedByMe = data.likes.some((user) => user._id === currentUserId);
  if (isLikedByMe) {
    likeButton.classList.add("card__like-button_is-active");
  }
  if (data.owner._id !== currentUserId) {
    deleteButton.remove();
  }
  if (onLikeIcon) {
    likeButton.addEventListener("click", () => onLikeIcon(likeButton, data._id, likeCountElement));
  }
  ...
  return cardElement;
};
\end{minted}

\subsection{Программный контроль ввода и валидация форм}
Модуль валидации обрабатывает данные форм на основе единого конфигурационного
объекта. Ошибки отображаются в реальном времени при вводе символов. Ключевая
логика опирается на объект состояния валидности поля.Обнаружение флага
patternMismatch означает нарушение регулярного выражения\cite{mdn_validity_state}.

\begin{minted}{javascript}
const checkInputValidity =
  (settings, inputElement, formElement) => {
    if (inputElement.validity.patternMismatch) {
      inputElement.setCustomValidity(
        inputElement.dataset.errorMessage
      );
    } else {
      inputElement.setCustomValidity("");
    }
    if (!inputElement.validity.valid) {
      showInputError(
        settings,
        inputElement,
        formElement,
        inputElement.validationMessage
      );
    } else {
      hideInputError(settings, inputElement, formElement);
    }
  };
\end{minted}

Блокировка кнопки подтверждения осуществляется путем обхода всех полей активной
формы. Функция очистки вызывается перед каждым открытием модального окна,
гарантируя возврат визуальных элементов к исходному чистому состоянию.

\subsection{Инициализация системы и обработчики отправки данных}
При загрузке страницы данные профиля и актуальный список публикаций
запрашиваются параллельно с помощью метода Promise.all. Данный метод
оптимизирует процесс инициализации интерфейса. Разрешение объединенного объекта
запускает процесс сохранения идентификатора пользователя, обновления текстовых
блоков и рендеринга сетки карточек.

\begin{minted}{javascript}
Promise.all([getUserInfo(), getCardList()])
  .then(([userData, cards]) => {
    profileTitle.textContent = userData.name;
    profileDescription.textContent = userData.about;
    profileAvatar.style.backgroundImage = `url(${userData.avatar})`;
    currentUserId = userData._id;
    cards.forEach((cardData) => {
      placesWrap.append(
        createCardElement(
          cardData,
          {
            onPreviewPicture: handlePreviewPicture,
            onLikeIcon: handleLikeCard,
            onDeleteCard: handleDeleteCard,
          },
          currentUserId
        )
      );
    });
  })
  .catch((err) => {
    console.log(`Ошибка при загрузке данных: ${err}`);
  });
\end{minted}

Обработчики форм следуют единому паттерну предотвращения стандартной отправки и
отображения статуса загрузки. Исходный текстовый контент управляющего элемента
кэшируется в отдельную локальную переменную. Заданное слово подставляется в
кнопку на время выполнения запроса. Блок окончательного выполнения finally
гарантирует стопроцентный возврат изначального текста после завершения фоновой
транзакции.

\begin{minted}{javascript}
const handleProfileFormSubmit = (evt) => {
  evt.preventDefault();
  const submitButton = evt.submitter;
  const initialText = submitButton.textContent;
  submitButton.textContent = 'Сохранение...';

  setUserInfo({
    name: profileTitleInput.value,
    about: profileDescriptionInput.value,
  })
    .then((userData) => {
      profileTitle.textContent = userData.name;
      profileDescription.textContent = userData.about;
      closeModalWindow(profileFormModalWindow);
    })
    .catch((err) => {
      console.log(`Ошибка при обновлении профиля: ${err}`);
    })
    .finally(() => {
      submitButton.textContent = initialText;
    });
};
\end{minted}

\subsection{Алгоритм агрегации и вывода статистики пользователей}
Дополнительное индивидуальное задание включает разработку системы сбора
аналитики. Событие клика по логотипу веб-приложения инициирует новый запрос к
серверу для получения максимально свежего массива публикаций. Алгоритм
осуществляет предварительную сортировку полученных карточек по дате создания.
Сортировка позволяет мгновенно извлечь даты первой и последней активности.
Временные метки конвертируются в локализованный строковый формат.

Для выявления абсолютного количества уникальных авторов и подсчета созданных
ими публикаций применяется специализированная коллекция Map. Идентификатор
пользователя выступает уникальным ключом словаря.

\begin{minted}{javascript}
const handleLogoClick = () => {
  getCardList()
    .then((cards) => {
      if (!cards || cards.length === 0) return;
      statsDefinitionsList.innerHTML = '';
      statsUsersList.innerHTML = '';
      statsTitle.textContent = 'Статистика пользователей';
      statsText.textContent = 'Все пользователи:';
      const totalCards = cards.length;
      const sortedCards = [...cards].sort((a, b) => new Date(a.createdAt) - new Date(b.createdAt));
      const firstCardDate = formatDate(new Date(sortedCards[0].createdAt));
      const lastCardDate = formatDate(new Date(sortedCards[sortedCards.length - 1].createdAt));
      const usersMap = new Map();
      sortedCards.forEach((card) => {
        const ownerId = card.owner._id;
        if (!usersMap.has(ownerId)) {
          usersMap.set(ownerId, { profile: card.owner, cardsCount: 0 });
        }
        usersMap.get(ownerId).cardsCount++;
      });
\end{minted}

Коллекция конвертируется в классический массив для проведения математических
вычислений. Выявление максимума публикаций на одного человека осуществляется
посредством глобального объекта Math.

Отображение агрегированных показателей производится через механизмы безопасной
шаблонизации. Специализированные вспомогательные функции клонируют скрытые узлы
разметки и наполняют их текстовой информацией. Готовые фрагменты добавляются в
очищенные контейнеры модального окна. Завершающим этапом алгоритма выступает
вызов функции переключения видимости аналитического интерфейса.

\begin{minted}{javascript}
      const uniqueUsersInfo = Array.from(usersMap.values());
      const totalUsers = uniqueUsersInfo.length;
      const maxCardsFromOneUser = Math.max(...uniqueUsersInfo.map(u => u.cardsCount));
      statsDefinitionsList.append(createInfoString("Всего карточек:", totalCards));
      statsDefinitionsList.append(createInfoString("Первая создана:", firstCardDate));
      statsDefinitionsList.append(createInfoString("Последняя создана:", lastCardDate));
      statsDefinitionsList.append(createInfoString("Всего пользователей:", totalUsers));
      statsDefinitionsList.append(createInfoString("Максимум карточек от одного:", maxCardsFromOneUser));
      uniqueUsersInfo.forEach((u) => {
        statsUsersList.append(createUserPreview(u.profile));
      });
      openModalWindow(statsModalWindow);
    })
    .catch((err) => {
      console.log(`Ошибка при загрузке статистики: ${err}`);
    });
};
\end{minted}

Разработанный механизм подтверждает технологическую возможность клиентских
скриптов самостоятельно выполнять сложные вычислительные операции над
неструктурированными массивами данных.

\conclusion
В ходе выполнения учебной практики был успешно спроектирован и разработан
полноценный клиентский веб-сервис «Mesto». Главная цель работы заключалась в
углублении теоретических знаний и закреплении практических навыков модульного
программирования на языке JavaScript с применением современных инструментов
автоматизации. Важным инфраструктурным достижением проекта стало развертывание
процесса непрерывной интеграции и доставки кодовой базы. Настройка
автоматизированных сценариев на платформе GitHub Actions обеспечила
автоматическую компиляцию исходного кода средствами сборщика Vite и последующее
размещение оптимизированного бандла на статическом хостинге GitHub Pages при
каждом обновлении репозитория \cite{github_actions, vite_docs}.

Генерация визуальных элементов интерфейса построена на базе безопасного
глубокого клонирования скрытых шаблонов разметки. Данный метод полностью
исключает риск возникновения уязвимостей межсайтового скриптинга при выводе
пользовательского контента. Процесс отображения графики оптимизирован за счет
минимизации прямых обращений к активному дереву узлов документа.

Формы ввода надежно защищены от отправки некорректной или неполной
информации. Разработанный скрипт в реальном времени анализирует содержимое
полей, управляет выводом кастомных текстовых подсказок и жестко блокирует
элементы управления до полного исправления всех допущенных ошибок \cite{mdn_validity_state}.

Успешно выполнено задание по разработке аналитического модуля. Модуль собирает
актуальный массив публикаций, осуществляет его предварительную хронологическую
сортировку и агрегирует статистику активности авторов. Применение
специализированной структуры данных коллекции Map позволило эффективно
вычислить абсолютные максимумы и сформировать списки уникальных участников
платформы с последующим выводом информации в модальное окно.

Проделанная работа подтверждает успешное освоение ключевых технологий
фронтенд-разработки, методологии автоматизации и принципов управления потоками
данных. Полученный практический опыт формирует надежный фундамент для
дальнейшего изучения компонентного мышления и проектирования сложных
распределенных информационных систем.

% Отобразить все источники. Даже те, на которые нет ссылок.

\bibliographystyle{ugost2003}
\bibliography{thesis}

% Окончание основного документа и начало приложений Каждая последующая секция
% документа будет являться приложением
\appendix

\section{Flash"=носитель с исходным кодом программ, использующихся в работе}
Исходный код приложен на съемном носителе и расположен в папке \texttt{mesto-ad}.

\begin{itemize}
    \item[\textbf{Папка}] \texttt{.github/workflows} "--- конфигурация GitHub Actions для автоматического деплоя
    \item[\textbf{Папка}] \texttt{src/blocks} "--- CSS"=файлы блоков интерфейса
    \item[\textbf{Папка}] \texttt{src/images} "--- графические ресурсы приложения
    \item[\textbf{Папка}] \texttt{src/vendor} "--- шрифты и нормализация стилей
    \item[\textbf{Папка}] \texttt{src/pages} "--- точка сборки стилей
    \item[\textbf{Папка}] \texttt{src/scripts} "--- js скрипты приложения
    \item[\textbf{Папка}] \texttt{src/scripts/components} "--- js модули приложения
    \item[\textbf{Файл}] \texttt{src/scripts/index.js} "--- точка входа приложения
    \item[\textbf{Файл}] \texttt{src/pages/index.css} "--- главный файл стилей
    \item[\textbf{Файл}] \texttt{.gitignore}
    \item[\textbf{Файл}] \texttt{README.md}
    \item[\textbf{Файл}] \texttt{index.html} "--- корневая HTML"=разметка приложения
    \item[\textbf{Файл}] \texttt{package.json} "--- зависимости и скрипты проекта
    \item[\textbf{Файл}] \texttt{vite.config.js} "--- конфигурация сборщика Vite
\end{itemize}

\section{Полный код api.js}
\begin{minted}{JavaScript}
// Конфигурация для запросов
const config = {
  baseUrl: "https://mesto.nomoreparties.co/v1/apf-cohort-203",
  headers: {
    authorization: "5478f85e-a2d0-4991-88ba-dd6636dc46d8",
    "Content-Type": "application/json",
  },
};

// Функция проверки ответа от сервера
const getResponseData = (res) => {
  return res.ok ? res.json() : Promise.reject(`Ошибка: ${res.status}`);
};

// Запрос на получение данных пользователя
export const getUserInfo = () => {
  return fetch(`${config.baseUrl}/users/me`, {
    headers: config.headers,
  }).then(getResponseData);
};

// Запрос на получение массива карточек
export const getCardList = () => {
  return fetch(`${config.baseUrl}/cards`, {
    headers: config.headers,
  }).then(getResponseData);
};

// Обновление данных профиля
export const setUserInfo = ({ name, about }) => {
  return fetch(`${config.baseUrl}/users/me`, {
    method: "PATCH",
    headers: config.headers,
    body: JSON.stringify({
      name,
      about,
    }),
  }).then(getResponseData);
};

// Обновление аватара пользователя
export const updateAvatar = (avatar) => {
  return fetch(`${config.baseUrl}/users/me/avatar`, {
    method: "PATCH",
    headers: config.headers,
    body: JSON.stringify({
      avatar,
    }),
  }).then(getResponseData);
};

// Добавление новой карточки
export const postNewCard = ({ name, link }) => {
  return fetch(`${config.baseUrl}/cards`, {
    method: "POST",
    headers: config.headers,
    body: JSON.stringify({
      name,
      link,
    }),
  }).then(getResponseData);
};

// Удаление карточки
export const deleteCardApi = (cardId) => {
  return fetch(`${config.baseUrl}/cards/${cardId}`, {
    method: "DELETE",
    headers: config.headers,
  }).then(getResponseData);
};

// Постановка и снятие лайка
export const changeLikeCardStatus = (cardId, isLiked) => {
  return fetch(`${config.baseUrl}/cards/likes/${cardId}`, {
    method: isLiked ? "DELETE" : "PUT",
    headers: config.headers,
  }).then(getResponseData);
};
\end{minted}

\end{document}

